\mode*
\subsection{Tools}

The development tool a student uses shapes which errors they meet and how they
recover from them. The most visible channel is the error message. Compiler and
interpreter messages are a well-documented obstacle for novices:
\textcite{Prather2017CEM} describe them as
\blockcquote{Prather2017CEM}{a major impediment to novice student
learning}
and a sizeable body of work has tried to rewrite them into something novices can
act on. Their human-factors study is cautiously encouraging --- after enhancing
the messages,
\blockcquote{Prather2017CEM}{most students are reading enhanced compiler error
messages and generally make effective changes after encountering them.}

This is the redeeming feature of the syntactic errors that dominate novices'
early mistakes: a modern IDE or compiler can catch them and point the student at
a fix, which is why \textcite{MisconceptionsSurvey2017} set them aside in favour
of conceptual difficulties:
\blockcquote{MisconceptionsSurvey2017}{problems in syntactic knowledge are often
easy to detect and fix. Perhaps that is why they are often noted as the most
frequent mistakes novices make [\ldots] A compiler or a modern integrated
development environment (IDE) may be able to find them and then provide error
messages or hints for correction.}

Whether a particular IDE should be recommended for CS1 is beyond our scope; the
point for us is that tooling can absorb much of the syntactic load, leaving the
conceptual misconceptions this paper addresses as the harder, more interesting
target.
